\section{归一化与 Dropout 改变训练时的随机系统}
\label{sec:14-Deep-Learning/03-Optimization-and-Training/03}

验证集准确率 \(94.3\%\)，损失曲线平滑收敛到 0.21，权重文件完整，预处理脚本和训练时同一份，随机种子也固定了。上线之后第一条真实样本进来，模型给出的概率是 0.03——同一条样本在验证集上给出过 0.87。

排查清单走完一轮找不到问题，是因为清单默认了一件事：训练时跑的和推理时跑的是同一个函数。有归一化层和 Dropout 的网络不满足这个前提。它在训练态定义的是一个随机映射，输入不只是当前样本，还包括同 batch 的其他样本和一次掩码采样；推理态把这些随机性关掉，换上一个确定性的替身。两个函数在同一个输入上给出多大差距，取决于替身做得有多准。

这两类层还有一层更基本的作用。它们不是用来纠正梯度的，而是用来改写梯度所在的那张曲面：归一化改变损失曲面的条件数和梯度对权重尺度的响应，Dropout 改变的是训练时被优化的目标本身。深度模型把选择表示交给优化器之后，人能施加的最有效影响，就落在这类改写曲面的操作上——它们没有对应的收敛定理，只有一批可复现的经验规律。

\subsection{一个样本的输出取决于同 batch 里的陌生人}

按直觉，某层在第三个样本上的激活 \(x_3=3.7\) 经过该层变换后的输出，应当只由 \(x_3\) 和该层参数决定。批归一化直接推翻这条直觉。对一个 mini-batch 中某通道的 \(m\) 个激活，它先算批统计量
\[
\mu_B=\frac1m\sum_{i=1}^m x_i,\qquad
\sigma_B^2=\frac1m\sum_{i=1}^m(x_i-\mu_B)^2,
\]
再输出
\[
y_i=\gamma\,\frac{x_i-\mu_B}{\sqrt{\sigma_B^2+\varepsilon}}+\beta,
\]
其中 \(\gamma,\beta\) 可学习，\(\varepsilon\) 防除零。

\(y_3\) 的表达式里带着 \(\mu_B\) 和 \(\sigma_B\)，而这两个量属于整个 batch。同一张图片，和一批亮度均值 0.6 的图片同批时被减掉 0.6，和一批均值 0.2 的图片同批时被减掉 0.2，除数也各不相同。训练过程中 batch 每个 epoch 重新洗牌，这个扰动一直在。

扰动不全是坏事。基准在动，网络就没法把结论压在某个激活的精确数值上，只能学一个在多种偏移和尺度下都成立的表示，这是一种由采样带来的正则。但代价同样确定：训练全程网络没有见过``输出只依赖当前样本''的那个确定性函数，在训练态的视角下每个样本的输出都是随机变量，方差来自同批的陌生样本。

推理时 batch 不存在，标准做法是换上训练期间累积的滑动统计量：
\[
\widehat\mu\leftarrow\alpha\widehat\mu+(1-\alpha)\mu_B,\qquad
\widehat\sigma^2\leftarrow\alpha\widehat\sigma^2+(1-\alpha)\sigma_B^2,
\]
\(\alpha\) 常取 0.9 或 0.99。至此推理函数成为确定性映射，但它是训练那个随机映射的一个近似替身，不是它的等价形式。

\subsection{归一化改的是曲面的条件数与梯度对尺度的响应}

批归一化能加速训练，通常被归因于``减少内部协变量偏移''，这个解释在后来的对照实验里并不稳固。更可靠的说法要看它对损失曲面做了什么。

先看一条可以证清楚的性质。设某层的权重矩阵为 \(W\)，其输出立即送进批归一化，且该层不带偏置。把 \(W\) 换成 \(aW\)（\(a>0\)）：预激活整体乘以 \(a\)，同一 batch 的均值与标准差也整体乘以 \(a\)，归一化后的结果完全不变，于是损失满足 \(L(aW)=L(W)\)。两边对 \(W\) 求微分，链式法则给出
\[
a\,\nabla L\big|_{aW}=\nabla L\big|_{W},
\qquad\text{即}\qquad
\norm{\nabla L\big|_{aW}}=\frac1a\norm{\nabla L\big|_{W}}.
\]
条件对账：要求归一化紧跟在该层之后、该层无偏置、\(\sigma_B>0\)、\(a>0\)，并且 \(\varepsilon\) 相对 \(\sigma_B^2\) 可忽略（否则等式退化为近似）。结论是权重范数与梯度范数成反比，于是权重涨大一倍，有效步长自动减半——权重范数在这里扮演了一个自动的学习率调节器，也解释了为什么权重衰减在带归一化的网络里会通过压小范数来间接抬高有效学习率。

再看条件数。归一化把每个通道的输出拉到大致相同的尺度，等价于对该层的输入做了一次逐通道的白化，损失曲面上原本因为各方向量级悬殊而拉长的等高线因此变圆。

\begin{center}
\begin{tikzpicture}[>=stealth]
\begin{scope}
\foreach \s in {0.45,0.72,1.0} {
  \draw[gray] (0,0) ellipse ({2.4*\s} and {0.5*\s});
}
\fill (0,0) circle (1.6pt);
\draw[very thick,->] (-2.2,0.42) -- (-1.70,-0.34) -- (-1.26,0.30)
  -- (-0.86,-0.24) -- (-0.52,0.19) -- (-0.24,-0.13) -- (0,0);
\node[below,font=\small] at (0,-1.35) {归一化前：条件数大};
\end{scope}
\begin{scope}[shift={(6.4,0)}]
\foreach \s in {0.45,0.72,1.0} {
  \draw[gray] (0,0) circle ({1.55*\s});
}
\fill (0,0) circle (1.6pt);
\draw[very thick,->] (-1.42,0.60) -- (0,0);
\node[below,font=\small] at (0,-1.35) {归一化后：条件数接近 1};
\end{scope}
\end{tikzpicture}
\end{center}

\emph{等高线越接近圆，负梯度方向越接近指向极小点，同一个学习率在各个方向上就都够用。}

图里左右两条轨迹用的是同一个优化器和同一个学习率，差别只在曲面。左边的窄谷里，负梯度几乎垂直于通往极小点的方向，步长被最陡的那一维卡死；右边的负梯度直指中心，一步顶左边六步。这也是归一化允许把学习率调高一个量级的原因：它没有让优化器变聪明，只是把优化器面对的问题换了一个。曲率与步长上界之间的关系见\hyperref[sec:09-Convex-Optimization/06-Optimization-Algorithms/01]{``梯度下降真正难在步长''}。

\subsection{批统计量的精度由 batch size 定，分布式训练会悄悄降级}

滑动统计量能不能代表真实分布，取决于每个 batch 拿多少样本去估计。单卡 batch size 为 4 时，通道均值的标准误差是 \(\sigma_{\text{true}}/\sqrt4=\sigma_{\text{true}}/2\)，每个 batch 给出的 \(\mu_B\) 在真值两侧大幅摆动。滑动平均能压低这个抖动，压不掉它——训练结束时 \(\widehat\mu\) 里残留的方差仍然可观，而推理只用这一个数。

分布式训练里同样的问题换了个外衣。八张卡各取 32 个样本，梯度经 all-reduce 平均，对优化器而言等效 batch size 是 256。但如果归一化的 \(\mu_B,\sigma_B^2\) 只在卡内计算、不跨卡同步，每张卡上的归一化仍然只看见 32 个样本。优化器按 256 的噪声水平设步长，归一化按 32 的噪声水平做变换，两边对同一次迭代的假设不一致。跨卡同步的归一化实现存在，但它要多一轮通信；不用它就得承认有效批统计量比名义 batch size 小得多。

\subsection{LayerNorm 用一条不含 batch 的轴换掉这个 gap}

层归一化对每个样本自己的特征维求统计量：
\[
\mu_{\mathrm{LN}}=\frac1D\sum_{j=1}^D x_j,\qquad
\sigma_{\mathrm{LN}}^2=\frac1D\sum_{j=1}^D(x_j-\mu_{\mathrm{LN}})^2.
\]
输出只由这个样本决定，训练态与推理态用的是同一个公式，前面那道 gap 直接消失。batch size 被迫很小甚至为 1 的场景——序列模型、大模型微调、在线单条推理——这一条往往就决定了选型。

代价在别处。跨特征维归一化会抹掉特征之间的绝对尺度差异，而这些差异有时正是信号：一层里某些单元惯常强激活、另一些惯常弱激活，相对大小携带信息。层归一化把它们统一拉到零均值单位方差，后续只能靠 \(\gamma,\beta\) 恢复一个全局的尺度与偏移，恢复不了单个样本上单个特征的特异幅度。选择在张量的哪条轴上做归一化，等于选择你愿意保留哪一类尺度信息、放弃哪一类，以及愿不愿意承担训练与推理的统计量不一致。

\subsection{Dropout 在训练时定义了一族模型，推理时只用一个}

Dropout 的随机源不同，造出的 gap 结构相同。设保留概率为 \(q\)，训练时对每个激活独立采样掩码 \(m_j\sim\Bernoulli(q)\)，倒置写法计算
\[
\widetilde h_j=\frac{m_j}{q}h_j,
\]
一阶期望被校准：\(\E[\widetilde h_j\given h_j]=h_j\)。推理时掩码关闭、不做缩放，因为训练期的 \(1/q\) 已经把期望对齐。

对齐只到第一个非线性为止。除非 \(f\) 是仿射的，否则
\[
\E[f(\widetilde h)]\neq f\bigl(\E[\widetilde h]\bigr).
\]
\(f\) 凸时 Jensen 不等式给出固定方向的偏差；ReLU 与 GELU 在整条实轴上既非凸也非凹，偏差的方向随激活分布而变，逐层累积之后没有简单的符号可言。含 \(n\) 个 Dropout 单元的网络在训练时定义的是 \(2^n\) 个子网络构成的一族模型，损失是这一族模型的期望损失；推理时那个确定性网络只是这族输出的一个近似代表，一阶匹配不蕴含分布匹配。

\begin{center}
\begin{tikzpicture}[x=1.05cm,y=1.0cm,>=stealth]
\foreach \x/\h in {0.3/0.15, 0.9/0.50, 1.5/1.05, 2.1/1.50, 2.7/1.35,
                   3.3/1.00, 3.9/0.70, 4.5/0.45, 5.1/0.25, 5.7/0.10} {
  \fill[gray!35] ({\x-0.26},0) rectangle ({\x+0.26},\h);
}
\draw[->] (-0.1,0) -- (6.3,0) node[right,font=\small] {输出};
\draw[dashed,very thick] (2.62,-0.15) -- (2.62,1.85);
\node[above,font=\scriptsize] at (2.62,1.85) {\(\E[f_{\text{train}}]\)};
\draw[very thick] (3.55,-0.15) -- (3.55,1.85);
\node[above,font=\scriptsize] at (3.75,1.85) {\(f_{\text{eval}}\)};
\node[below,font=\scriptsize] at (3.0,-0.25) {各掩码子网络的输出分布};
\end{tikzpicture}
\end{center}

\emph{关掉掩码得到的那个确定性输出，落在子网络输出分布里的位置由非线性决定，未必对准分布的均值。}

图里那道实线与虚线之间的距离，就是训练与推理之间的功能偏移。偏移小的时候一切正常，这也是 Dropout 大多数时候能用的原因；偏移大的时候，验证指标依然漂亮，因为验证通常也走推理态，两边错得一致，只有真实数据分布稍有变化时才会暴露。

关于强度还有一条边界。随机丢弃迫使信息在单元之间分散，共适应被打断，这是它的正则效应；但当模型容量本就不足、数据本就稀缺，或者归一化层已经在注入批噪声时，再叠一层高强度 Dropout，损失的是拟合能力而不是过拟合风险。近年的视觉与语言主干里 Dropout 比例普遍下调甚至取消，正是因为归一化和数据增强已经把噪声预算占掉了。

\subsection{训练态与推理态是两张计算图}

把前面几段收拢：带这两类层的网络同时定义了 \(f_{\text{train}}(x,\text{batch},\text{mask})\) 与 \(f_{\text{eval}}(x)\)，后者不是前者的平均，而是去掉两个噪声源之后的一个近似。由此产生的故障有一个共同特征——不报错。

保存权重时漏掉归一化层的 \(\widehat\mu,\widehat\sigma^2\)，加载后它们是零或随机初值，每个通道的归一化基准整体错位，网络照常输出数值。微调时把 batch size 设成 1 或 2 却继续更新滑动统计量，预训练阶段积累的良好估计被极噪声的批统计量冲掉，推理用的是一份被污染的统计量。验证时忘记切到 eval 模式，Dropout 仍在丢弃单元，验证曲线被无谓的噪声抬高，好模型可能因此被判为差模型。归一化与 Dropout 的前后顺序不同，两态之间的统计差异量级也不同，公开结论时有互相矛盾，部分原因就是各实验在这一点上没有对齐。数值层面的静默失效——inf 与 NaN 的首次出现位置、对数域中的稳定写法——见\hyperref[ch:067]{``机器学习中的稳定计算''}。

所以是否使用这两类层，不该由``深层网络都用''来决定，而该由几个可以事先回答的问题决定：训练初期的 loss 是否震荡到需要归一化来压条件数；有效 batch 是否大到支撑得起可信的批统计量，撑不起就换一条不含 batch 的归一化轴；部署时是否逐样本推理；训练数据与部署数据之间是否存在协变量偏移，因为滑动统计量在偏移下会系统性地失准。这些都是训练与推理一致性的工程决策，正则化强度的调参排在它们后面。

到这里，一个能跑起来并且能被诊断的训练系统就搭齐了：步长与数值窗口、带记忆的更新、被改写过的曲面。剩下的问题是，这样训练出来的网络到底学到了什么——参数表里那几亿个数字并不直接回答这个问题，本部分最后一个单元处理它。
